% Generated draft: requires mathematical and source review.
\documentclass[11pt]{article}
\usepackage[margin=25mm]{geometry}
\usepackage{fontspec}
\setmainfont{Latin Modern Roman}
\usepackage{amsmath,amssymb,mathtools,unicode-math}
\setmathfont{Latin Modern Math}
\usepackage{enumitem,needspace,xurl,hyperref}
\setlength{\parindent}{0pt}
\setlength{\parskip}{5pt}
\setlength{\emergencystretch}{3em}
\newcommand{\N}{\mathbb{N}}
\newcommand{\Z}{\mathbb{Z}}
\newcommand{\Q}{\mathbb{Q}}
\newcommand{\R}{\mathbb{R}}
\newcommand{\C}{\mathbb{C}}
\newcommand{\F}{\mathbb{F}}
\newcommand{\A}{\mathbb{A}}
\newcommand{\PP}{\mathbb{P}}
\newcommand{\E}{\mathbb{E}}
\newcommand{\eps}{\varepsilon}
\newcommand{\dd}{\,\mathrm{d}}
\begin{document}
\section*{UnsolvedMath ID: 2100206}
\textbf{Source problem code:} catalogue entry
\subsection{Definitions and mathematical statement}
An integrable Hamiltonian system with $n$ degrees of freedom has a symplectic phase space $(M,\omega)$ of dimension $2n$ and a map $F=(f_1,\ldots,f_n)$ whose components Poisson commute and have independent differentials on a dense regular set. A point is singular if $\operatorname{rank}\dd F<n$; a fiber is singular if it contains such a point.

The problem asks, for $1\leq n\leq3$, to describe all topological types of singularities possible in \emph{algebraically integrable} systems, and to describe their symplectic invariants. The supplied statement does not define the algebraic category. In particular it does not specify real versus complex systems, regular versus meromorphic maps, properness or compactification, or whether generic fibers must be abelian varieties or open subsets of them. Those are not silently added to the elementary integrability framework above.

The classification also requires an equivalence convention. Topology of a point germ, topology of a whole singular fiber, and topology of the fibration near that fiber are distinct data. A symplectic classification can similarly concern germs or neighborhoods of entire fibers, and can either fix $F$ or permit reparametrization of its base. No such convention is provided. Thus the task remains a source-dependent low-dimensional classification program, rather than a uniquely specified list of singularities or invariants.

The supplied assessment claims a one-degree-of-freedom classification by $A_1$ points (real elliptic or hyperbolic forms) and $A_2$ cusps with Milnor number $\mu\leq2$, and claims completeness of an action germ with characteristic behaviors $h$, $h\log|h|$, and $h^{5/6}$. Its needed fiber/genus/irreducibility hypotheses and precise action/energy conventions are not given; in particular reduction of a fiber alone is not a supplied justification for that classification. The assessment itself says the $h^{5/6}$ scaling was not fully verified analytically. These statements are retained as unverified supplied progress claims, not promoted to a general theorem in the source's unspecified category.
\subsection{Short English statement}
Classify the topology and symplectic invariants of singularities in algebraically integrable systems with at most three degrees of freedom. The supplied assessment labels the task partially solved, claims a one-degree result with the qualifications above, and says the two- and three-degree cases remain open. It mentions degenerations of abelian varieties, genus-two/Jacobian classifications, semi-global symplectic invariants and stable degenerate singularities. It also stresses continuous symplectic moduli for degenerate singularities, referring to source Problem 2.3, and reports no general degenerate analogue of Eliasson's theorem. These are attributed survey claims rather than independent proofs or a definitive current literature review.
\subsection{Sources}
\begin{itemize}
\item UnsolvedMath permanent ID 2100206, source code catalogue entry, ``Open Problems in Integrable Systems --- Topology of integrable systems, Lagrangian fibrations, and their invariants.'' Registry curation and metadata: Ulam AI and UnsolvedMath Contributors, CC BY 4.0; updated 2026-08-06.
\item \emph{Open Problems, Questions, and Challenges in Finite-Dimensional Integrable Systems} (2018), Problem 2.6: \url{https://arxiv.org/abs/1804.03737}. The supplied assessment also points to Problem 2.3. Source-TeX extraction accessed 2026-07-29; status/release and rights review are NEEDS\_REVIEW.
\item The supplied summary names Arinkin--Fedorov, Namikawa--Ueno, Vũ Ngọc, Kalashnikov, and Eliasson, but gives no complete bibliographic entries or theorem locators. No missing reference or author list for the source volume is inferred.
\end{itemize}
\end{document}
